% Options for packages loaded elsewhere
\PassOptionsToPackage{unicode}{hyperref}
\PassOptionsToPackage{hyphens}{url}
%
\documentclass[
]{article}
\usepackage{amsmath,amssymb}
\usepackage{iftex}
\ifPDFTeX
  \usepackage[T1]{fontenc}
  \usepackage[utf8]{inputenc}
  \usepackage{textcomp} % provide euro and other symbols
\else % if luatex or xetex
  \usepackage{unicode-math} % this also loads fontspec
  \defaultfontfeatures{Scale=MatchLowercase}
  \defaultfontfeatures[\rmfamily]{Ligatures=TeX,Scale=1}
\fi
\usepackage{lmodern}
\ifPDFTeX\else
  % xetex/luatex font selection
\fi
% Use upquote if available, for straight quotes in verbatim environments
\IfFileExists{upquote.sty}{\usepackage{upquote}}{}
\IfFileExists{microtype.sty}{% use microtype if available
  \usepackage[]{microtype}
  \UseMicrotypeSet[protrusion]{basicmath} % disable protrusion for tt fonts
}{}
\makeatletter
\@ifundefined{KOMAClassName}{% if non-KOMA class
  \IfFileExists{parskip.sty}{%
    \usepackage{parskip}
  }{% else
    \setlength{\parindent}{0pt}
    \setlength{\parskip}{6pt plus 2pt minus 1pt}}
}{% if KOMA class
  \KOMAoptions{parskip=half}}
\makeatother
\usepackage{xcolor}
\usepackage{color}
\usepackage{fancyvrb}
\newcommand{\VerbBar}{|}
\newcommand{\VERB}{\Verb[commandchars=\\\{\}]}
\DefineVerbatimEnvironment{Highlighting}{Verbatim}{commandchars=\\\{\}}
% Add ',fontsize=\small' for more characters per line
\newenvironment{Shaded}{}{}
\newcommand{\AlertTok}[1]{\textcolor[rgb]{1.00,0.00,0.00}{\textbf{#1}}}
\newcommand{\AnnotationTok}[1]{\textcolor[rgb]{0.38,0.63,0.69}{\textbf{\textit{#1}}}}
\newcommand{\AttributeTok}[1]{\textcolor[rgb]{0.49,0.56,0.16}{#1}}
\newcommand{\BaseNTok}[1]{\textcolor[rgb]{0.25,0.63,0.44}{#1}}
\newcommand{\BuiltInTok}[1]{\textcolor[rgb]{0.00,0.50,0.00}{#1}}
\newcommand{\CharTok}[1]{\textcolor[rgb]{0.25,0.44,0.63}{#1}}
\newcommand{\CommentTok}[1]{\textcolor[rgb]{0.38,0.63,0.69}{\textit{#1}}}
\newcommand{\CommentVarTok}[1]{\textcolor[rgb]{0.38,0.63,0.69}{\textbf{\textit{#1}}}}
\newcommand{\ConstantTok}[1]{\textcolor[rgb]{0.53,0.00,0.00}{#1}}
\newcommand{\ControlFlowTok}[1]{\textcolor[rgb]{0.00,0.44,0.13}{\textbf{#1}}}
\newcommand{\DataTypeTok}[1]{\textcolor[rgb]{0.56,0.13,0.00}{#1}}
\newcommand{\DecValTok}[1]{\textcolor[rgb]{0.25,0.63,0.44}{#1}}
\newcommand{\DocumentationTok}[1]{\textcolor[rgb]{0.73,0.13,0.13}{\textit{#1}}}
\newcommand{\ErrorTok}[1]{\textcolor[rgb]{1.00,0.00,0.00}{\textbf{#1}}}
\newcommand{\ExtensionTok}[1]{#1}
\newcommand{\FloatTok}[1]{\textcolor[rgb]{0.25,0.63,0.44}{#1}}
\newcommand{\FunctionTok}[1]{\textcolor[rgb]{0.02,0.16,0.49}{#1}}
\newcommand{\ImportTok}[1]{\textcolor[rgb]{0.00,0.50,0.00}{\textbf{#1}}}
\newcommand{\InformationTok}[1]{\textcolor[rgb]{0.38,0.63,0.69}{\textbf{\textit{#1}}}}
\newcommand{\KeywordTok}[1]{\textcolor[rgb]{0.00,0.44,0.13}{\textbf{#1}}}
\newcommand{\NormalTok}[1]{#1}
\newcommand{\OperatorTok}[1]{\textcolor[rgb]{0.40,0.40,0.40}{#1}}
\newcommand{\OtherTok}[1]{\textcolor[rgb]{0.00,0.44,0.13}{#1}}
\newcommand{\PreprocessorTok}[1]{\textcolor[rgb]{0.74,0.48,0.00}{#1}}
\newcommand{\RegionMarkerTok}[1]{#1}
\newcommand{\SpecialCharTok}[1]{\textcolor[rgb]{0.25,0.44,0.63}{#1}}
\newcommand{\SpecialStringTok}[1]{\textcolor[rgb]{0.73,0.40,0.53}{#1}}
\newcommand{\StringTok}[1]{\textcolor[rgb]{0.25,0.44,0.63}{#1}}
\newcommand{\VariableTok}[1]{\textcolor[rgb]{0.10,0.09,0.49}{#1}}
\newcommand{\VerbatimStringTok}[1]{\textcolor[rgb]{0.25,0.44,0.63}{#1}}
\newcommand{\WarningTok}[1]{\textcolor[rgb]{0.38,0.63,0.69}{\textbf{\textit{#1}}}}
\usepackage{longtable,booktabs,array}
\usepackage{calc} % for calculating minipage widths
% Correct order of tables after \paragraph or \subparagraph
\usepackage{etoolbox}
\makeatletter
\patchcmd\longtable{\par}{\if@noskipsec\mbox{}\fi\par}{}{}
\makeatother
% Allow footnotes in longtable head/foot
\IfFileExists{footnotehyper.sty}{\usepackage{footnotehyper}}{\usepackage{footnote}}
\makesavenoteenv{longtable}
\setlength{\emergencystretch}{3em} % prevent overfull lines
\providecommand{\tightlist}{%
  \setlength{\itemsep}{0pt}\setlength{\parskip}{0pt}}
\setcounter{secnumdepth}{-\maxdimen} % remove section numbering
\ifLuaTeX
  \usepackage{selnolig}  % disable illegal ligatures
\fi
\usepackage{bookmark}
\IfFileExists{xurl.sty}{\usepackage{xurl}}{} % add URL line breaks if available
\urlstyle{same}
\hypersetup{
  pdftitle={Dark Matter as the Metric Shadow of Dead History},
  hidelinks,
  pdfcreator={LaTeX via pandoc}}

\title{Dark Matter as the Metric Shadow of Dead History}
\author{}
\date{}

\begin{document}
\maketitle

{
\setcounter{tocdepth}{3}
\tableofcontents
}
\section{Dark Matter as the Metric Shadow of Dead
History}\label{dark-matter-as-the-metric-shadow-of-dead-history}

\subsection{A Flux/MIP Archive Model for Galactic Rotation
Curves}\label{a-fluxmip-archive-model-for-galactic-rotation-curves}

\textbf{Author:} Jamie McCabe, with AI-assisted formalization\\
\textbf{Version:} v3q paper bundle, 2026-05-12

\begin{center}\rule{0.5\linewidth}{0.5pt}\end{center}

\subsection{Abstract}\label{abstract}

The standard interpretation of galactic rotation curves introduces a
non-luminous collisionless mass component, while modified-gravity
alternatives usually encode the observed excess acceleration through a
universal acceleration scale. This paper develops a different
interpretation within the Flux/MIP framework: galactic dark-sector
behaviour is treated as the metric response to historical registration
structure. In this view, baryonic matter is not only a source of
mass-energy but also a persistent archive of prior actualization events.
We define a global spatial archive depth, introduce a saturating
archival-efficiency function, and extend the model with a
compact-remnant archive amplifier motivated by white dwarfs, neutron
stars, and stellar/rogue black holes as high historical-depth nodes. The
resulting v3q model is tested against the official SPARC175 mass-model
catalogue. In the present implementation, Newtonian baryons alone give a
galaxy-mean MAE of approximately \textbf{39.98 km/s}, the v3o spatial
archive model gives \textbf{27.96 km/s}, and the v3q compact-remnant
archive model gives \textbf{18.71 km/s}. The model does not yet solve
SPARC at MOND-level precision, but it substantially improves over the
baryonic baseline and isolates a structured residual pathway associated
with old, gas-poor, remnant-rich systems. The central thesis is:
\textbf{dark matter is the metric shadow of dead history.}

\begin{center}\rule{0.5\linewidth}{0.5pt}\end{center}

\subsection{1. Introduction}\label{introduction}

The galactic rotation-curve problem is usually stated as a missing-mass
problem. Observed circular velocities remain too high at large radii
compared with the Newtonian prediction from visible baryons. The
standard cosmological answer is cold dark matter: an extended
collisionless halo surrounding galaxies. However, the tight empirical
coupling between baryons and acceleration suggests that the dark-sector
signal is not arbitrary. It tracks baryonic structure with remarkable
fidelity.

The Flux/MIP programme reframes this coupling. Instead of asking where
unseen mass is located, we ask how the metric stores, exports, and
responds to historical registration. In this framework, matter is
interpreted as bound history: a stable residue of measurement-induced
actualization events. Mass is therefore not merely substance in space,
but persistent causal record.

The present paper develops the v3q version of this idea for galaxy
rotation curves. The model starts with a global archive-depth scalar
derived from baryonic surface compression, then adds a saturating
efficiency law and a compact-remnant wake amplifier. The SPARC175
dataset provides a first falsification benchmark.

\begin{center}\rule{0.5\linewidth}{0.5pt}\end{center}

\subsection{2. Flux/MIP Ontology}\label{fluxmip-ontology}

The Flux/MIP framework treats physical reality as a system of
actualization events. MIP has two associated meanings:

\begin{enumerate}
\def\labelenumi{\arabic{enumi}.}
\tightlist
\item
  \textbf{Measurement-Induced Projection:} the transition from
  unresolved quantum potential into classical registration.
\item
  \textbf{Minimum Intersection Point:} the local condition under which a
  possibility becomes a persistent event in the metric.
\end{enumerate}

Matter is the stable residue of repeated actualization. The nucleus of
an atom is interpreted as a hardened archive of prior registrations,
while electron shells form the active interface through which the object
exchanges information and entropy with its environment.

The galactic model therefore distinguishes between:

\begin{itemize}
\tightlist
\item
  \textbf{spatial compression}, the present baryonic density state;
\item
  \textbf{temporal persistence}, the accumulated duration of stable
  registration;
\item
  \textbf{compact-remnant texture}, the hidden population of high-depth
  archive nodes.
\end{itemize}

\begin{center}\rule{0.5\linewidth}{0.5pt}\end{center}

\subsection{3. Mass as Bound History}\label{mass-as-bound-history}

In standard dynamics, inertial mass is a scalar resistance to
acceleration. In the Flux/MIP interpretation, inertial mass is the work
required to re-register an already persistent historical record. This
gives a compact statement:

\begin{quote}
Mass is persistent bound history.
\end{quote}

Gravity is then the metric response to gradients and concentrations of
this bound historical structure.

This interpretation does not require changing local inertial mass
directly. The model instead modifies the effective nonlocal wake
produced by the galaxy-level archive. This avoids immediately violating
equivalence-principle constraints, because the correction is treated as
a metric response rather than a composition-dependent inertial-mass
change.

\begin{center}\rule{0.5\linewidth}{0.5pt}\end{center}

\subsection{4. Black Holes as Horizon-Locked Historical
Archives}\label{black-holes-as-horizon-locked-historical-archives}

Black holes are treated as the limiting case of historical registration.
A collapse event occurs when accumulated temporal depth and coherence
exceed the local entropy-export capacity:

\begin{Shaded}
\begin{Highlighting}[]
\NormalTok{D\_t C\_sigma / S\_leak \textgreater{} Theta\_BH}
\end{Highlighting}
\end{Shaded}

where:

\begin{itemize}
\tightlist
\item
  \texttt{D\_t} is temporal depth,
\item
  \texttt{C\_sigma} is coherence/compression,
\item
  \texttt{S\_leak} is entropy export capacity,
\item
  \texttt{Theta\_BH} is the horizon-lock threshold.
\end{itemize}

The event horizon is interpreted as a causal vault: a topological
boundary that quarantines saturated historical registration from the
surrounding metric. In this language, black holes are not singularities
of matter but singularities of accumulated history.

This motivates v3q. Mature galaxies contain compact remnants: white
dwarfs, neutron stars, stellar black holes, and potentially rogue black
holes. They need not dominate baryonic mass, but they may dominate
historical-depth texture.

\begin{center}\rule{0.5\linewidth}{0.5pt}\end{center}

\subsection{5. The v3q Master Equations}\label{the-v3q-master-equations}

\subsubsection{5.1 Global Spatial Archive
Depth}\label{global-spatial-archive-depth}

The galaxy is assigned a global baryonic archive depth:

\begin{Shaded}
\begin{Highlighting}[]
\NormalTok{D\_spatial = (Sigma\_bar / Sigma\_ref)\^{}0.08}
\end{Highlighting}
\end{Shaded}

with:

\begin{Shaded}
\begin{Highlighting}[]
\NormalTok{Sigma\_bar = M\_bar / (2 pi R\_disk\^{}2)}
\NormalTok{M\_bar = M\_star + M\_gas}
\NormalTok{M\_star = Upsilon\_star L\_3.6}
\NormalTok{M\_gas = 1.33 M\_HI}
\NormalTok{Sigma\_ref = 4.08e9 Msun kpc\^{}{-}2}
\NormalTok{Upsilon\_star = 0.5}
\end{Highlighting}
\end{Shaded}

\subsubsection{5.2 Saturating Archival
Efficiency}\label{saturating-archival-efficiency}

The v3o causal-sluice efficiency is:

\begin{Shaded}
\begin{Highlighting}[]
\NormalTok{eta(D) = 8.0 D\^{}{-}7 / [1 + (D / 0.68)\^{}8]}
\end{Highlighting}
\end{Shaded}

This term prevents runaway archive leakage at high depth.

\subsubsection{5.3 Memory Radius}\label{memory-radius}

The radial expression scale of the wake is:

\begin{Shaded}
\begin{Highlighting}[]
\NormalTok{L\_chi = 7.9 R\_disk D\^{}{-}4.6}
\end{Highlighting}
\end{Shaded}

Low-depth systems generate broader wakes; high-depth systems compress
their archive more tightly.

\subsubsection{5.4 Universal Wake Texture}\label{universal-wake-texture}

The enclosed historical wake mass is:

\begin{Shaded}
\begin{Highlighting}[]
\NormalTok{M\_hist(\textless{}r) = eta\_eff M\_bar x\^{}s / (1+x)\^{}(s{-}1)}
\end{Highlighting}
\end{Shaded}

with:

\begin{Shaded}
\begin{Highlighting}[]
\NormalTok{x = r / L\_chi}
\NormalTok{s = 1.36}
\end{Highlighting}
\end{Shaded}

\subsubsection{5.5 Compact-Remnant Archive
Proxy}\label{compact-remnant-archive-proxy}

The raw remnant proxy is:

\begin{Shaded}
\begin{Highlighting}[]
\NormalTok{R\_rem = (1 {-} f\_gas) ((10 {-} T) / 10)}
\NormalTok{f\_gas = M\_gas / M\_bar}
\end{Highlighting}
\end{Shaded}

where \texttt{T} is Hubble type. The term targets old, gas-poor systems.

\subsubsection{5.6 Retention Gate}\label{retention-gate}

A physically conservative variant requires a sufficiently deep potential
to retain and organize compact archive nodes:

\begin{Shaded}
\begin{Highlighting}[]
\NormalTok{G\_ret = (V\_max / V\_ref)\^{}n / [1 + (V\_max / V\_ref)\^{}n]}
\end{Highlighting}
\end{Shaded}

A useful physical gate explored here is:

\begin{Shaded}
\begin{Highlighting}[]
\NormalTok{V\_ref = 200 km/s}
\NormalTok{n = 4}
\end{Highlighting}
\end{Shaded}

The retained remnant archive is:

\begin{Shaded}
\begin{Highlighting}[]
\NormalTok{R\_eff = R\_rem G\_ret}
\end{Highlighting}
\end{Shaded}

\subsubsection{5.7 Effective Archival
Efficiency}\label{effective-archival-efficiency}

The v3q remnant amplifier is:

\begin{Shaded}
\begin{Highlighting}[]
\NormalTok{eta\_eff = eta(D) [1 + chi\_rem R\_eff\^{}zeta]}
\end{Highlighting}
\end{Shaded}

with the empirical best ungated values:

\begin{Shaded}
\begin{Highlighting}[]
\NormalTok{chi\_rem = 3.0}
\NormalTok{zeta = 0.5}
\end{Highlighting}
\end{Shaded}

The square-root dependence implies sublinear amplification:

\begin{Shaded}
\begin{Highlighting}[]
\NormalTok{eta\_eff = eta(D) [1 + 3 sqrt(R\_eff)]}
\end{Highlighting}
\end{Shaded}

Compact remnants therefore act as a coherence amplifier, not as a linear
mass reservoir.

\subsubsection{5.8 Total Rotation Curve}\label{total-rotation-curve}

The historical acceleration is:

\begin{Shaded}
\begin{Highlighting}[]
\NormalTok{g\_hist(r) = G M\_hist(\textless{}r) / r\^{}2}
\end{Highlighting}
\end{Shaded}

The baryonic velocity contribution is computed from SPARC mass-model
components:

\begin{Shaded}
\begin{Highlighting}[]
\NormalTok{V\_bar\^{}2 = |V\_gas| V\_gas + Upsilon\_disk V\_disk\^{}2 + Upsilon\_bulge V\_bul\^{}2}
\end{Highlighting}
\end{Shaded}

with:

\begin{Shaded}
\begin{Highlighting}[]
\NormalTok{Upsilon\_disk = 0.5}
\NormalTok{Upsilon\_bulge = 0.7}
\end{Highlighting}
\end{Shaded}

The final prediction is:

\begin{Shaded}
\begin{Highlighting}[]
\NormalTok{V\_Flux\^{}2(r) = V\_bar\^{}2(r) + G eta\_eff M\_bar / r * x\^{}1.36 / (1+x)\^{}0.36}
\end{Highlighting}
\end{Shaded}

\begin{center}\rule{0.5\linewidth}{0.5pt}\end{center}

\subsection{6. SPARC175 Validation}\label{sparc175-validation}

The model was evaluated against the SPARC175 rotation-curve dataset
using the official global galaxy table and mass-model table. The
implementation parses the global table for \texttt{L\_3.6},
\texttt{R\_disk}, \texttt{M\_HI}, \texttt{T}, and quality values, and
parses the mass-model table for \texttt{R}, \texttt{Vobs},
\texttt{Vgas}, \texttt{Vdisk}, and \texttt{Vbul}.

\subsubsection{6.1 Summary Metrics}\label{summary-metrics}

\begin{longtable}[]{@{}
  >{\raggedright\arraybackslash}p{(\columnwidth - 6\tabcolsep) * \real{0.2143}}
  >{\raggedleft\arraybackslash}p{(\columnwidth - 6\tabcolsep) * \real{0.2857}}
  >{\raggedleft\arraybackslash}p{(\columnwidth - 6\tabcolsep) * \real{0.2857}}
  >{\raggedright\arraybackslash}p{(\columnwidth - 6\tabcolsep) * \real{0.2143}}@{}}
\toprule\noalign{}
\begin{minipage}[b]{\linewidth}\raggedright
Model
\end{minipage} & \begin{minipage}[b]{\linewidth}\raggedleft
Point MAE
\end{minipage} & \begin{minipage}[b]{\linewidth}\raggedleft
Galaxy Mean MAE
\end{minipage} & \begin{minipage}[b]{\linewidth}\raggedright
Notes
\end{minipage} \\
\midrule\noalign{}
\endhead
\bottomrule\noalign{}
\endlastfoot
Newtonian baryons & 47.01 km/s & 39.98 km/s & Fixed baryonic
mass-to-light ratios \\
v3o spatial archive & 33.01 km/s & 27.96 km/s & Causal-sluice only \\
v3p-T morphology proxy & 28.76 km/s & 27.96 km/s & Helpful point-level
correction, too blunt galaxy-level \\
v3q compact-remnant archive & 19.59 km/s & 18.71 km/s & Best empirical
branch \\
v3q-retention scan best & 19.75 km/s & 18.92 km/s & Conservative
retention-gated branch \\
\end{longtable}

\subsubsection{6.2 Strongly Improved
Systems}\label{strongly-improved-systems}

\begin{longtable}[]{@{}lrrrrr@{}}
\toprule\noalign{}
Galaxy & T & f\_gas & D\_spatial & v3o MAE & v3q MAE \\
\midrule\noalign{}
\endhead
\bottomrule\noalign{}
\endlastfoot
UGC02487 & 0 & 0.010 & 0.861 & 127.98 & 44.71 \\
NGC6674 & 3 & 0.038 & 0.844 & 79.80 & 15.44 \\
NGC3992 & 4 & 0.019 & 0.873 & 73.03 & 16.42 \\
UGC02885 & 5 & 0.026 & 0.801 & 74.77 & 19.45 \\
NGC5985 & 3 & 0.015 & 0.820 & 110.22 & 59.31 \\
NGC2841 & 3 & 0.014 & 0.903 & 101.94 & 53.55 \\
NGC5907 & 5 & 0.031 & 0.846 & 54.44 & 6.75 \\
UGC12506 & 6 & 0.063 & 0.791 & 84.76 & 41.29 \\
\end{longtable}

\subsubsection{6.3 Overcorrection Cases}\label{overcorrection-cases}

The first v3q proxy over-amplifies some intermediate systems. These
cases motivate the retention-gated form and/or a more refined
remnant-population proxy.

\begin{longtable}[]{@{}lrrrrr@{}}
\toprule\noalign{}
Galaxy & T & f\_gas & D\_spatial & v3o MAE & v3q MAE \\
\midrule\noalign{}
\endhead
\bottomrule\noalign{}
\endlastfoot
NGC4088 & 4 & 0.020 & 0.913 & 14.19 & 41.07 \\
NGC6195 & 3 & 0.014 & 0.773 & 12.36 & 38.11 \\
NGC4051 & 4 & 0.007 & 0.822 & 6.47 & 31.81 \\
NGC4389 & 4 & 0.007 & 0.792 & 29.91 & 51.52 \\
NGC5371 & 4 & 0.009 & 0.845 & 15.95 & 36.41 \\
UGC11557 & 8 & 0.054 & 0.761 & 13.51 & 28.33 \\
NGC4217 & 3 & 0.008 & 0.877 & 36.04 & 48.55 \\
F561-1 & 9 & 0.096 & 0.699 & 5.62 & 15.54 \\
\end{longtable}

\begin{center}\rule{0.5\linewidth}{0.5pt}\end{center}

\subsection{7. Interpretation}\label{interpretation}

The v3q result suggests that the residual signal in massive early-type
systems is not random. The systems that improve most are old, gas-poor,
high-luminosity galaxies where compact-remnant archive density should be
highest. The square-root scaling indicates that the first compact
archive nodes matter disproportionately; after a remnant population
exists, additional remnants yield diminishing returns.

This is distinct from MACHO dark matter. The model does not claim that
stellar remnants supply all the missing gravitational mass. Instead,
compact remnants are treated as high-registration-depth nodes that
amplify a nonlocal metric wake.

\begin{center}\rule{0.5\linewidth}{0.5pt}\end{center}

\subsection{8. Failure Modes and Falsification
Tests}\label{failure-modes-and-falsification-tests}

A serious model must be falsifiable. The following tests are direct:

\begin{enumerate}
\def\labelenumi{\arabic{enumi}.}
\tightlist
\item
  \textbf{Residual-remnant correlation:} v3q predicts that v3o residuals
  should correlate with old stellar populations, low gas fraction, and
  compact-remnant indicators.
\item
  \textbf{Color and stellar-population tests:} External SDSS/WISE/GALEX
  colours should reduce scatter if temporal persistence is real.
\item
  \textbf{Lensing consistency:} Any archive wake that affects rotation
  curves must also have a consistent weak-lensing signature.
\item
  \textbf{Bullet Cluster-style systems:} If baryonic archive sources and
  lensing peaks are fully separated without a historical-lag mechanism,
  the model fails.
\item
  \textbf{Solar System screening:} The archive wake must vanish in
  high-coherence local environments.
\item
  \textbf{Independent SPARC reruns:} The model must survive quality
  cuts, alternate mass-to-light ratios, and independent parsers.
\end{enumerate}

\begin{center}\rule{0.5\linewidth}{0.5pt}\end{center}

\subsection{9. Limitations}\label{limitations}

The current v3q model is not a final replacement for ΛCDM or MOND. It is
a structured prototype. Its galaxy-mean MAE remains above MOND-level
performance, and the remnant proxy is coarse. Hubble type and gas
fraction are not direct measurements of compact-remnant distributions.
The retention gate is physically motivated but not yet derived from
first principles.

Therefore, the correct empirical claim is modest:

\begin{quote}
v3q substantially improves over Newtonian baryons and identifies a
physically structured residual pathway tied to compact-remnant archive
physics.
\end{quote}

\begin{center}\rule{0.5\linewidth}{0.5pt}\end{center}

\subsection{10. Conclusion}\label{conclusion}

Flux/MIP v3q converts the dark matter question from a missing-substance
problem into a missing-history problem. A galaxy is not merely a
distribution of visible mass. It is a historical archive containing
smooth baryonic depth, temporal persistence, and compact remnant
texture.

The present SPARC175 implementation shows a clear improvement from
Newtonian baryons to v3o and from v3o to v3q:

\begin{Shaded}
\begin{Highlighting}[]
\NormalTok{39.98 km/s {-}\textgreater{} 27.96 km/s {-}\textgreater{} 18.71 km/s}
\end{Highlighting}
\end{Shaded}

for galaxy-mean MAE.

The strongest conceptual result is:

\begin{quote}
Dark-sector behaviour is not missing matter alone; it is missing
historical registration structure.
\end{quote}

Or, in the compact thesis statement:

\begin{quote}
\textbf{Dark matter is the metric shadow of dead history.}
\end{quote}

\begin{center}\rule{0.5\linewidth}{0.5pt}\end{center}

\subsection{References}\label{references}

\begin{itemize}
\tightlist
\item
  Lelli, F., McGaugh, S. S., \& Schombert, J. M. (2016). SPARC: Mass
  Models for 175 Disk Galaxies with Spitzer Photometry and Accurate
  Rotation Curves. \emph{Astronomical Journal}.
\item
  McGaugh, S. S., Lelli, F., \& Schombert, J. M. (2016). The Radial
  Acceleration Relation in Rotationally Supported Galaxies.
  \emph{Physical Review Letters}.
\item
  Milgrom, M. (1983). A modification of the Newtonian dynamics as a
  possible alternative to the hidden mass hypothesis.
  \emph{Astrophysical Journal}.
\item
  Bekenstein, J. D. (1973). Black holes and entropy. \emph{Physical
  Review D}.
\item
  Hawking, S. W. (1975). Particle creation by black holes.
  \emph{Communications in Mathematical Physics}.
\end{itemize}

\end{document}
